\section{例题与分析}

\begin{question}\begin{example}[表头改装的两种量程]\label{c6:example:1}
表头满偏电流 $I_g$、内阻 $R_g$。分别改成量程 $U$ 的电压表和量程 $I>I_g$ 的电流表，求所需电阻、仪表总内阻及附加电阻满量程时的功率。
\end{example}
\typeout{ECFIG:fig:c6:example:1:s}\fig{fig:c6:example:1:s}{表头的满偏电流方向及对应端电压}{\begin{tikzpicture}

\draw(0,0)to[ammeter,l={$I_g,R_g$}](4,0);
\draw[flow](.2,1)--(1.5,1)node[midway,above]{$I_g$};
\node[left]at(0,0){$+$};\node[right]at(4,0){$-$};

\end{tikzpicture}}
\end{question}
\begin{solution}\textbf{先由表头的满偏条件求电压量程}\par\nobreak

图\ref{fig:c6:example:1:s}取电流流入端为电压正端，表头满偏压降为 $U_g=I_gR_g$。要扩大量程须有 $U>U_g$；串联电阻承担剩余压降，且与表头通过同一电流。
\[
 R_s=\frac{U-I_gR_g}{I_g}=\frac U{I_g}-R_g,\qquad
 R_V=R_s+R_g=\frac U{I_g}.
\]

\typeout{ECFIG:fig:c6:example:1:p}\fig{fig:c6:example:1:p}{上图串联分压，下图并联分流；两条支路电流由同一端口电压约束}{\begin{tikzpicture}

\draw(0,2)to[R,l={$R_s$}](3,2)to[ammeter,l={$R_g$}](6,2);
\draw[flow](.2,3.1)--(1.8,3.1)node[midway,above]{$I_g$};
\draw[<->](0,1.1)--(6,1.1)node[midway,below]{$U=I_g(R_s+R_g)$};
\begin{scope}[yshift=-1.2cm]
\draw(-1,-1)--(0,-1)--(0,0)to[ammeter,l={$R_g$}](6,0)--(6,-1)--(7,-1);
\draw(0,-1)--(0,-2.5)to[R,l_={$R_p$}](6,-2.5)--(6,-1);
\fill(0,-1)circle(2pt);\fill(6,-1)circle(2pt);
\draw[flow](-1,-.4)--(-.1,-.4)node[midway,above]{$I$};
\draw[flow](.4,-.6)--(1.8,-.6)node[midway,below]{$I_g$};
\draw[flow](.4,-1.8)--(1.8,-1.8)node[midway,above]{$I-I_g$};
\end{scope}

\end{tikzpicture}}


\par\medskip\Needspace{4\baselineskip}\noindent\textbf{用并联同电压确定分流电阻}\par\nobreak

图\ref{fig:c6:example:1:p}下图中，端口输入 $I$，表头只能通过 $I_g$，故其余 $I-I_g$ 必须经 $R_p$。由 $I_gR_g=(I-I_g)R_p$ 得
\[
 R_p=\frac{I_gR_g}{I-I_g},\qquad R_A=\frac{I_gR_g}{I}.
\]


\par\medskip\Needspace{4\baselineskip}\noindent\textbf{分别计算附加电阻的额定发热}\par\nobreak

满量程时串联电阻功率 $P_s=I_g^2R_s$，分流电阻功率
\[
P_p=(I-I_g)^2R_p=I_gR_g(I-I_g).
\]
两功率都对应附加电阻自身，未把表头的 $I_g^2R_g$ 算入。扩大量程时均须核对电阻的功率承受能力；分流电阻阻值虽小，流过的电流却可能较大。
\end{solution}
\begin{question}\begin{example}[半偏法的仪器选择与改装]\label{c6:example:2}
表头 $I_g=500\,\mu\mathrm A$、内阻约 $100\,\Omega$。有 $999.9\,\Omega$、$99999.9\,\Omega$ 两电阻箱，以及 $200\,\Omega$、$2\,\mathrm{k}\Omega$ 滑动变阻器和 $2,6\,\mathrm V$ 电源。用半偏法测内阻，选择串联与并联调节元件；若测得 $R_g=105.0\,\Omega$，改成 $10\,\mathrm V$ 表应串联多大电阻？
\end{example}
\typeout{ECFIG:fig:c6:example:2:s}\fig{fig:c6:example:2:s}{半偏实验先断开并联开关，再保持串联电阻不变合上开关}{\begin{tikzpicture}

\draw(0,0)to[battery1,invert,l={$E$}](0,3.5)to[R,l={$R_1$}](3.5,3.5)to[ammeter,l={$G$}](3.5,0)--(0,0);
\draw(3.5,3.5)to[nos,l={$S$}](6,3.5)to[R,l={$R_2$}](6,0)--(3.5,0);
\fill(3.5,3.5)circle(2pt);\fill(3.5,0)circle(2pt);

\end{tikzpicture}}
\end{question}
\begin{solution}\textbf{由满偏所需串联阻值选择电源和电阻箱}\par\nobreak

图\ref{fig:c6:example:2:s}先断开 $S$，调 $R_1$ 使表头满偏。用 $6\,\mathrm V$ 电源时
\[
R_1\simeq\frac{6}{500\times10^{-6}}-100=11900\,\Omega.
\]
因此串联采用 $99999.9\,\Omega$ 电阻箱。若用 $2\,\mathrm V$，所需 $3900\,\Omega$ 也须大箱，但 $R_g/R_1$ 较大；在给定器材中用 $6\,\mathrm V$ 更能减小半偏近似误差。


\par\medskip\Needspace{4\baselineskip}\noindent\textbf{保持串联阻值不变再调并联分流}\par\nobreak

合上 $S$ 后只调 $R_2$，使表头减为 $250\,\mu\mathrm A$。所需 $R_2$ 接近 $100\,\Omega$，采用 $999.9\,\Omega$ 电阻箱便于准确读数。两个滑动变阻器可作粗调，但不能代替并联支路的准确阻值读数。

为检验近似，满偏时$E=I_g(R_1+R_g)$；半偏时并联端电压为$I_gR_g/2$，分流电流为$I_gR_g/(2R_2)$，故
\[
E=R_1\left(\frac{I_g}{2}+\frac{I_gR_g}{2R_2}\right)+\frac{I_gR_g}{2}.
\]
消去$E,I_g$，先得$2(R_1+R_g)=R_1+R_1R_g/R_2+R_g$，再得
\[
R_2(R_1+R_g)=R_1R_g,\qquad
R_2=\frac{R_1R_g}{R_1+R_g}<R_g.
\]
相同电源下并联后总电流略增，因而不能把并联电阻与表头电阻精确等同。


\par\medskip\Needspace{4\baselineskip}\noindent\textbf{把测得的表头内阻代入扩大量程条件}\par\nobreak

\typeout{ECFIG:fig:c6:example:2:p}\fig{fig:c6:example:2:p}{改装后的整个端口承受十伏；表头本身只承受 $I_gR_g$}{\begin{tikzpicture}

\draw(0,0)to[R,l={$19895\,\Omega$}](3.5,0)to[ammeter,l={$105\,\Omega$}](7,0);
\draw[flow](.5,1.1)--(2,1.1)node[midway,above]{$500\,\mu\mathrm A$};
\draw[<->](0,-1)--(7,-1)node[midway,below]{$10\,\mathrm V$};

\end{tikzpicture}}

按题给 $R_g=105.0\,\Omega$ 作为已测得内阻，图\ref{fig:c6:example:2:p}要求
\[
R_s=\frac{10}{500\times10^{-6}}-105=19895\,\Omega,
\qquad R_V=20000\,\Omega.
\]
检查：附加电阻压降 $9.9475\,\mathrm V$，表头压降 $0.0525\,\mathrm V$，合计恰为 $10\,\mathrm V$。若 $105.0\,\Omega$ 只是未经修正的半偏箱读数，应先修正，不能混作真实内阻。
\end{solution}
\begin{question}\begin{example}[同一电压表依次测串联大电阻]\label{c6:example:3}
$R_1,R_2$ 串联接理想恒压源 $U$。内阻 $R_V$ 的同一电压表依次跨接两电阻，读数 $U_1,U_2$。求两读数、比值与和，判断有限内阻是否必然使比值也失真。
\end{example}
\typeout{ECFIG:fig:c6:example:3:s}\fig{fig:c6:example:3:s}{第一次测量：电压表跨在 $A,M$ 间，恒压源始终接在 $A,B$ 间}{\begin{tikzpicture}

\draw(0,0)to[battery1,invert,l={$U$}](0,3)to[R,l={$R_1$}](3,3)to[R,l={$R_2$}](6,3)--(6,0)--(0,0);
\draw(0,3)--(0,4.8)to[voltmeter,l={$R_V;\ U_1$}](3,4.8)--(3,3);
\foreach \x/\lab in {0/A,3/M,6/B}{\fill(\x,3)circle(2pt);\node[below]at(\x,3){$\lab$};}

\end{tikzpicture}}
\end{question}
\begin{solution}\textbf{把第一次测量的被测电阻与电压表并联}\par\nobreak

图\ref{fig:c6:example:3:s}中 $R_1'=R_1R_V/(R_1+R_V)$ 与 $R_2$ 串联，$A,B$ 间电压仍为 $U$。所以
\[
 U_1=U\frac{R_1'}{R_1'+R_2}
 =\frac{UR_1R_V}{D},\qquad D=R_1R_2+R_V(R_1+R_2).
\]


\par\medskip\Needspace{4\baselineskip}\noindent\textbf{换接电压表后重新建立第二个电路}\par\nobreak

\typeout{ECFIG:fig:c6:example:3:p}\fig{fig:c6:example:3:p}{第二次跨接 $M,B$；此时被分流的是 $R_2$，不与第一次测量同时发生}{\begin{tikzpicture}

\draw(0,0)to[battery1,invert,l={$U$}](0,3)to[R,l={$R_1$}](3,3)to[R,l={$R_2$}](6,3)--(6,0)--(0,0);
\draw(3,3)--(3,4.8)to[voltmeter,l={$R_V;\ U_2$}](6,4.8)--(6,3);
\foreach \x/\lab in {0/A,3/M,6/B}{\fill(\x,3)circle(2pt);\node[below]at(\x,3){$\lab$};}

\end{tikzpicture}}

图\ref{fig:c6:example:3:p}中改为 $R_2'=R_2R_V/(R_2+R_V)$，于是
\[
U_2=U\frac{R_2'}{R_1+R_2'}=\frac{UR_2R_V}{D}.
\]


\par\medskip\Needspace{4\baselineskip}\noindent\textbf{比较共同分母并解释读数之和}\par\nobreak

由于是同一只表，两式有共同分母，故 $U_1/U_2=R_1/R_2$ 精确成立；但
\[
U_1+U_2=U\frac{R_V(R_1+R_2)}{R_1R_2+R_V(R_1+R_2)}<U.
\]
两读数来自不同电路，不能当成同一时刻沿 $A\to M\to B$ 的两个电势降相加。$R_V\to\infty$ 时加载消失，读数之和才趋于 $U$。比值保留来自代数抵消，并不表示两个读数各自准确。
\end{solution}
\begin{question}\begin{example}[两种伏安接法的修正]\label{c6:example:4}
测得 $U_V=6.00\,\mathrm V,I_A=0.100\,\mathrm A$，已知 $R_A=1.00\,\Omega,R_V=3000\,\Omega$。分别求电压表仅跨 $R_x$ 与电压表跨 $R_x+R_A$ 时的真实电阻及未修正误差。
\end{example}
\typeout{ECFIG:fig:c6:example:4:s}\fig{fig:c6:example:4:s}{题设：两种伏安接法的修正}{\begin{tikzpicture}[scale=.88]
\draw(0,0)to[battery1,invert,l={$E$}](0,3)to[ammeter,l={$R_A$}](2,3)--(4,3)to[R,l={$R_x$}](4,0)--(0,0);\draw(2.6,3)--(6,3)to[voltmeter,l={$R_V$}](6,0)--(4,0);\node[above]at(4,3){$P$};\node[below]at(4,0){$Q$};
\end{tikzpicture}}
\end{question}
\begin{solution}\textbf{电压表只跨待测电阻时扣除分流}\par\nobreak

图\ref{fig:c6:example:4:s}中 $P,Q$ 间电压就是 $U_x$，电流表却读两支路电流之和。由节点 $P$ 的电流守恒，
\[
 I_V=\frac6{3000}=0.00200\,\mathrm A,\qquad
 I_x=I_A-I_V=0.0980\,\mathrm A,
\]
故 $R_x=6/0.0980=61.2245\,\Omega$。未修正值 $R_{\text{测}}=60.0\,\Omega$，相对误差为
\[
\frac{R_{\text{测}}-R_x}{R_x}=-2.00\%.
\]


\par\medskip\Needspace{4\baselineskip}\noindent\textbf{电压表跨电流表与待测电阻时扣除压降}\par\nobreak

\typeout{ECFIG:fig:c6:example:4:p}\fig{fig:c6:example:4:p}{电流表只读 $I_x$，电压表跨整个 $P,Q$ 端口，含 $R_A$ 的压降}{\begin{tikzpicture}

\draw(0,0)to[battery1,invert,l={$E$}](0,3)to[ammeter,l={$R_A$}](3,3)to[R,l={$R_x$}](3,0)--(0,0);
\draw(0,3)--(-2.5,3)to[voltmeter,l_={$R_V$}](-2.5,0)--(0,0);
\fill(0,3)circle(2pt)node[above]{$P$};\fill(0,0)circle(2pt)node[below]{$Q$};
\draw[flow](3.8,2.7)--(3.8,1.9)node[midway,right]{$I_A=I_x$};

\end{tikzpicture}}

图\ref{fig:c6:example:4:p}中 $I_x=I_A$，而 $U_x=U_V-I_AR_A$，所以
\[
 R_x=\frac{U_V}{I_A}-R_A=59.0\,\Omega,
 \qquad\frac{60.0-59.0}{59.0}=+1.69\%.
\]


\par\medskip\Needspace{4\baselineskip}\noindent\textbf{用接线位置检查两种误差方向}\par\nobreak

分流使所用电流偏大，商 $U/I$ 偏小；分压使所用电压偏大，商偏大。两种接法虽有相同读数，却对应不同的真实待测电阻，不能把两个修正结果当成同一电阻的重复测量。
\end{solution}
\begin{question}\begin{example}[欧姆表的中值与读数]\label{c6:example:5}
串联欧姆表电池 $E=1.5\,\mathrm V$，表头 $I_g=300\,\mu\mathrm A,R_g=100\,\Omega$，电池内阻 $0.5\,\Omega$。短接调满偏后，求调节电阻、中值电阻；测量时表头电流为满偏的 $1/3$，求待测电阻。
\end{example}
\typeout{ECFIG:fig:c6:example:5:s}\fig{fig:c6:example:5:s}{题设：欧姆表的中值与读数}{\begin{tikzpicture}[scale=.88]
\draw(0,0)to[battery1,l={$E,\ r_B=0.5\,\Omega$}](0,3)to[ammeter,l={$I_g,R_g$}](2.5,3)to[R,l={$R_0$}](5,3)--(5,0)to[R,l_={$R_x$}](0,0);\node[below]at(0,0){黑};\node[below]at(5,0){红};
\end{tikzpicture}}
\end{question}
\begin{solution}\textbf{用短接满偏确定内部总电阻}\par\nobreak
图\ref{fig:c6:example:5:s} 中短接时内部总电阻 $R_m=E/I_g=1.5/(300\times10^{-6})=5000\,\Omega$。除去表头和电池内阻，调节电阻 $R_0=5000-100-0.5=4899.5\,\Omega$。


\par\medskip\Needspace{4\baselineskip}\noindent\textbf{由非均匀刻度建立电流比关系}\par\nobreak
\typeout{ECFIG:fig:c6:example:5:p}\fig{fig:c6:example:5:p}{解答过程：欧姆表的中值与读数}{\begin{tikzpicture}[scale=.88]
\draw[world](160:2.2)arc(160:20:2.2);\draw(160:2.1)--(160:2.3);\node at(160:2.65){$\infty$};\draw(146.0:2.1)--(146.0:2.3);\node at(146.0:2.65){$45$};\draw(132.0:2.1)--(132.0:2.3);\node at(132.0:2.65){$20$};\draw(113.33333333333334:2.1)--(113.33333333333334:2.3);\node at(113.33333333333334:2.65){$10$};\draw(90.0:2.1)--(90.0:2.3);\node at(90.0:2.65){$5$};\draw(66.66666666666667:2.1)--(66.66666666666667:2.3);\node at(66.66666666666667:2.65){$2.5$};\draw(43.33333333333333:2.1)--(43.33333333333333:2.3);\node at(43.33333333333333:2.65){$1$};\draw(20:2.1)--(20:2.3);\node at(20:2.65){$0$};\draw[force](0,0)--(113.33333333333334:2);\fill(0,0)circle(2pt);\node at(0,-.6){$\Omega\times10^3$};
\end{tikzpicture}}

图\ref{fig:c6:example:5:p} 电流比 $f=I/I_g$ 满足 $f=R_m/(R_m+R_x)$，所以 $R_x=R_m(1/f-1)$。

\par\medskip\Needspace{4\baselineskip}\noindent\textbf{分别代入半偏与三分之一偏转并检查方向}\par\nobreak
半偏 $f=1/2$ 对应中值 $5.00\,\mathrm{k}\Omega$；三分之一偏转对应 $R_x=10.0\,\mathrm{k}\Omega$。电阻读数与电流不成正比，不能把三分之一偏转读成量程的三分之一。\end{solution}
\begin{question}\begin{example}[从低阻测量换到千欧测量]\label{c6:example:6}
模拟多用表刚用 $\Omega\times1$ 测过约 $10\,\Omega$ 电阻，接着测约 $2\,\mathrm{k}\Omega$。中值基础刻度为 $20\,\Omega$。选择倍率、调零顺序，并说明图示半偏时读数。
\end{example}
\typeout{ECFIG:fig:c6:example:6:s}\fig{fig:c6:example:6:s}{题设：从低阻测量换到千欧测量}{\begin{tikzpicture}[scale=.88]
\draw[world](160:2.2)arc(160:20:2.2);\draw(160:2.1)--(160:2.3);\node at(160:2.65){$\infty$};\draw(146.0:2.1)--(146.0:2.3);\node at(146.0:2.65){$180$};\draw(132.0:2.1)--(132.0:2.3);\node at(132.0:2.65){$80$};\draw(113.33333333333334:2.1)--(113.33333333333334:2.3);\node at(113.33333333333334:2.65){$40$};\draw(90.0:2.1)--(90.0:2.3);\node at(90.0:2.65){$20$};\draw(66.66666666666667:2.1)--(66.66666666666667:2.3);\node at(66.66666666666667:2.65){$10$};\draw(43.33333333333333:2.1)--(43.33333333333333:2.3);\node at(43.33333333333333:2.65){$4$};\draw(20:2.1)--(20:2.3);\node at(20:2.65){$0$};\draw[force](0,0)--(90.0:2);\fill(0,0)circle(2pt);\node at(0,-.6){$\Omega\quad\text{基础刻度}$};
\end{tikzpicture}}
\end{question}
\begin{solution}\textbf{让待测阻值落在中值附近}\par\nobreak

图\ref{fig:c6:example:6:s}基础中值为 $20\,\Omega$，若选择倍率 $k$，实际中值为 $20k\,\Omega$。测约 $2\,\mathrm{k}\Omega$，应取 $k=100$，使指针接近半偏、刻度较易分辨。


\par\medskip\Needspace{4\baselineskip}\noindent\textbf{先短接调零，再接待测电阻}\par\nobreak

\typeout{ECFIG:fig:c6:example:6:p}\fig{fig:c6:example:6:p}{同一基础刻度上的两个状态：短接时到右端零欧，接待测电阻后回到中值}{\begin{tikzpicture}

\draw[world](160:3)arc(160:20:3);
\foreach \a/\lab in {160/\infty,146/180,132/80,113.333/40,90/20,66.667/10,43.333/4,20/0}{\draw(\a:2.9)--(\a:3.1);\node at(\a:3.55){$\lab$};}
\draw[->,gray,thick](0,0)--(20:2.8);\draw[->,blue,thick](0,0)--(90:2.8);
\node[right]at(2.7,.2){短接调零};\node[left,blue]at(-.4,1.6){接 $R_x$};
\node at(0,-.7){$\Omega\times100$};\fill(0,0)circle(2pt);

\end{tikzpicture}}

改变倍率后，先短接两表笔，调欧姆零使指针到图\ref{fig:c6:example:6:p}右端 $0\,\Omega$；再断开表笔，接入与外部电源隔离的电阻。机械零点若原来正确，无须每次换倍率重调机械零点。


\par\medskip\Needspace{4\baselineskip}\noindent\textbf{把基础读数乘以倍率}\par\nobreak

接入后半偏对应基础刻度 $20$，所以 $R_x=20\times100=2.0\,\mathrm{k}\Omega$。在 $\times1$ 挡测同一电阻，$I/I_g=20/(20+2000)\simeq0.00990$，几乎停在无穷大端，说明换挡改善的是读数分辨率。
\end{solution}
\begin{question}\begin{example}[电池变弱后欧姆表的系统偏差]\label{c6:example:7}
原标定电池电动势 $E$，满偏电流 $I_g$。后来电池电动势降为 $E^\prime<E$，仍能通过调节电阻使表笔短接满偏。沿用原电阻刻度时，同一待测电阻 $R_x$ 将读大还是读小？求读数表达式。
\end{example}
\typeout{ECFIG:fig:c6:example:7:s}\fig{fig:c6:example:7:s}{题设：电池变弱后欧姆表的系统偏差}{\begin{tikzpicture}[scale=.88]
\draw(0,0)to[battery1,l={$E$}](0,3)to[ammeter,l={$I_g,R_g$}](2.5,3)to[R,l={$R_0$}](5,3)--(5,0)to[R,l_={$R_x$}](0,0);\node[below]at(0,0){黑};\node[below]at(5,0){红};
\end{tikzpicture}}
\end{question}
\begin{solution}\textbf{把重新调零后的内部电阻作为新中值}\par\nobreak

图\ref{fig:c6:example:7:s}原标定满足 $R_m=E/I_g$。电池变弱后仍能短接满偏，说明调节后的实际内部总电阻为 $R_m'=E'/I_g<R_m$。接入同一 $R_x$ 时
\[
 f=\frac{I}{I_g}=\frac{R_m'}{R_m'+R_x},\qquad
 R_x=R_m'\left(\frac1f-1\right).
\]

\typeout{ECFIG:fig:c6:example:7:p}\fig{fig:c6:example:7:p}{同一电流比按原刻度换算会读大；图仅以 $E^\prime/E=2/3$ 示意，两交点共用同一条竖直线}{\begin{tikzpicture}

\draw[axis](0,0)--(5.2,0)node[right]{$f$};\draw[axis](0,0)--(0,4)node[above]{$R/R_m$};
\draw[blue,thick,domain=.25:1,samples=80]plot({4.5*\x},{1/\x-1});
\draw[red,thick,domain=.25:1,samples=80]plot({4.5*\x},{(2/3)*(1/\x-1)});
\draw[guide](2.25,0)--(2.25,1);\fill[blue](2.25,1)circle(2pt);\fill[red](2.25,.666667)circle(2pt);
\draw[guide](2.25,1)--(3.3,2)node[right,blue]{$R_{\text{读}}/R_m$};
\draw[guide](2.25,.666667)--(3.3,.8)node[right,red]{$R_x/R_m$};
\node[below]at(4.5,0){$1$};\node[below]at(2.25,0){同一 $f$};

\end{tikzpicture}}


\par\medskip\Needspace{4\baselineskip}\noindent\textbf{按原刻度换算实际偏转}\par\nobreak

图\ref{fig:c6:example:7:p}原刻度仍使用旧中值 $R_m$，故
\[
 R_{\text{读}}=R_m(1/f-1)=\frac{R_m}{R_m'}R_x
 =\frac E{E'}R_x>R_x.
\]


\par\medskip\Needspace{4\baselineskip}\noindent\textbf{核对短路端与整条刻度的区别}\par\nobreak

$R_x=0$ 时两式都给零；但任意非零电阻的读数均被放大 $E/E'$ 倍。调零只保证短路端正确。电池内阻增大会限制可调零范围；只要仍可调零，上述串联模型的偏差因子由电动势变化决定。
\end{solution}
\begin{question}\begin{example}[模拟欧姆表测二极管反向电阻]\label{c6:example:8}
给定模拟串联欧姆表，黑表笔接内电池正侧，红表笔接负侧。如何连接二极管以测反向电阻？换量程后需怎样操作？为什么不同挡读数可能不同？
\end{example}
\typeout{ECFIG:fig:c6:example:8:s}\fig{fig:c6:example:8:s}{题设：模拟欧姆表测二极管反向电阻}{\begin{tikzpicture}[scale=.88]
\draw(0,0)to[battery1,l={$E$}](0,3)to[ammeter,l={$I_g,R_g$}](2.5,3)to[R,l={$R_0$}](5,3)--(5,0)to[R,l_={$R_x$}](0,0);\node[below]at(0,0){黑};\node[below]at(5,0){红};\node at(2.5,-1){被测物换成二极管};
\end{tikzpicture}}
\end{question}
\begin{solution}\textbf{由内电池极性判断二极管的反向接法}\par\nobreak

图\ref{fig:c6:example:8:s}的外部测量电流由黑表笔流出。反向偏置要求阴极电势高于阳极，所以黑笔接阴极、红笔接阳极。若相反连接，则变成正向测量。

\typeout{ECFIG:fig:c6:example:8:p}\fig{fig:c6:example:8:p}{二极管接入欧姆表的完整回路；竖线一侧为阴极，黑笔使其处于较高电势}{\begin{tikzpicture}

\draw(0,0)to[battery1,l={$E$}](0,3.5)to[ammeter,l={$R_g$}](3,3.5)to[R,l={$R_0$}](6,3.5)--(6,0)to[D](0,0);
\node[below]at(0,0){黑：阴极 $+$};\node[below]at(6,0){红：阳极 $-$};
\draw[flow](1,-1)--(2.5,-1)node[midway,below]{$I_{\text{反}}$};

\end{tikzpicture}}


\par\medskip\Needspace{4\baselineskip}\noindent\textbf{换挡后重做欧姆调零并检查实际偏转}\par\nobreak

换倍率后先短接表笔重新调零，再按图\ref{fig:c6:example:8:p}接回二极管。若偏转太小，读数分辨率有限；不能为追求中间刻度而任意提高未核对过的测试电压。


\par\medskip\Needspace{4\baselineskip}\noindent\textbf{把读数解释为该工作点的电压电流比}\par\nobreak

不同挡的内电源及限流电阻可能改变，二极管的反向电流也随实际端电压变化。因此仪表读出的 $U_D/I_{\text{反}}$ 属于该工作点，不能当作与挡位无关的恒定电阻。反向漏电很小时，指针可始终靠近无穷大端，这本身并不要求把它调到半偏。
\end{solution}
\begin{question}\begin{example}[用欧姆表观察电容充电与漏电]\label{c6:example:9}
一个起初未带电的电容接到串联欧姆表两笔之间。若理想无漏电，描述指针变化。若稳定后表头仍有满偏电流的 $1/6$，仪表中值电阻为 $R_m$，求等效漏电电阻。
\end{example}
\typeout{ECFIG:fig:c6:example:9:s}\fig{fig:c6:example:9:s}{题设：用欧姆表观察电容充电与漏电}{\begin{tikzpicture}[scale=.88]
\draw(0,0)to[battery1,l={$E$}](0,3)to[ammeter,l={$I_g,R_g$}](2.5,3)to[R,l={$R_0$}](5,3)--(5,0)to[C,l_={$C$}](0,0);\node[below]at(0,0){黑};\node[below]at(5,0){红};
\end{tikzpicture}}
\end{question}
\begin{solution}\textbf{先看接入瞬间与最终稳态}\par\nobreak

图\ref{fig:c6:example:9:s}的电容初电压为零。以 $R_m$ 表示调零后的内部总串联电阻，接入瞬间近似短路，$i(0^+)=E/R_m=I_g$；随后电容充电。记电容电压为$u(t)$，从电源到电容列回路方程$E=R_mi+u$，再用$i=C\,\dd u/\dd t$，得到
\[
R_mC\frac{\dd u}{\dd t}=E-u,\qquad u(0)=0.
\]
令$w=E-u$，则$\dd w=-\dd u$、$w(0)=E$，于是
\[
\frac{\dd w}{w}=-\frac{\dd t}{R_mC},\qquad
\int_E^{w(t)}\frac{\dd \xi}{\xi}
=-\int_0^t\frac{\dd s}{R_mC}.
\]
两边的原函数和端点分别为
\[
 [\ln\xi]_E^{w(t)}=-\frac1{R_mC}[s]_0^t,\qquad
 \ln\frac{w(t)}E=-\frac{t}{R_mC}.
\]
初值 $E>0$，偏差 $w$ 保持正号，取指数后
\[
w(t)=Ee^{-t/(R_mC)},\qquad
u(t)=E\bigl[1-e^{-t/(R_mC)}\bigr].
\]
代回$i=(E-u)/R_m=w/R_m$，得到
\[
i(t)=I_g\exp[-t/(R_mC)]\longrightarrow0.
\]
指针从小电阻侧向大电阻侧回移，最终指向无穷大。


\par\medskip\Needspace{4\baselineskip}\noindent\textbf{把稳定电流对应到电容的并联漏电通路}\par\nobreak

\typeout{ECFIG:fig:c6:example:9:p}\fig{fig:c6:example:9:p}{欧姆表内电源通过 $R_m$ 给 $C$ 与 $R_\ell$ 供电；稳态电流只走漏电支路}{\begin{tikzpicture}

\draw(0,0)to[battery1,invert,l={$E$}](0,3)to[R,l={$R_m$}](3,3)--(4,3)to[C,l={$C$}](4,0)--(0,0);
\draw(4,3)--(7,3)to[R,l={$R_\ell$}](7,0)--(4,0);
\fill(4,3)circle(2pt);\fill(4,0)circle(2pt);
\draw[flow](7.9,2.5)--(7.9,1.6)node[midway,right]{$I_\infty=I_g/6$};
\node at(4,-.9){稳态 $i_C=0$};

\end{tikzpicture}}

图\ref{fig:c6:example:9:p}中稳态电容支路不通直流，故
\[
\frac{I_\infty}{I_g}=\frac{E/(R_m+R_\ell)}{E/R_m}=\frac16,
\qquad R_\ell=5R_m.
\]


\par\medskip\Needspace{4\baselineskip}\noindent\textbf{核对暂态方向和测量端口}\par\nobreak

有漏电时，图\ref{fig:c6:example:9:p}的输入电流分成充电电流与漏电电流，故
\[
\frac{E-u}{R_m}=C\frac{\dd u}{\dd t}+\frac{u}{R_\ell}.
\]
移项后定义
\[
u_\infty=\frac{ER_\ell}{R_m+R_\ell},\qquad
\tau=\frac{CR_mR_\ell}{R_m+R_\ell},
\]
方程化为$\dd u/\dd t=(u_\infty-u)/\tau$。令$z=u_\infty-u$，则$\dd z/z=-\dd t/\tau$。初始未带电，$z(0)=u_\infty$，积分得
\[
\begin{aligned}
\int_{u_\infty}^{z(t)}\frac{\dd \xi}{\xi}&=-\int_0^t\frac{\dd s}{\tau},\\
[\ln\xi]_{u_\infty}^{z(t)}&=-\frac1\tau[s]_0^t,\\
\ln\frac{z(t)}{u_\infty}&=-\frac t\tau.
\end{aligned}
\]
于是$u(t)=u_\infty[1-e^{-t/\tau}]$，仪表读到的总电流为
\begin{align*}
i(t)&=\frac{E-u(t)}{R_m}
=\frac{E-u_\infty}{R_m}+\frac{u_\infty}{R_m}e^{-t/\tau}\\
&=I_\infty+(I_g-I_\infty)e^{-t/\tau}.
\end{align*}
这里$(E-u_\infty)/R_m=E/(R_m+R_\ell)=I_\infty$，且$\tau=(R_m\parallel R_\ell)C$。因此电流由$I_g$减小到非零的$I_\infty$，指针回移后停在有限阻值。若电容仍与其他并联元件连接，测得的是整个端口等效，不能把稳定电流全部归因于电容漏电。
\end{solution}
\begin{question}\begin{example}[六伏三瓦灯泡的伏安实验]\label{c6:example:10}
灯泡额定 $6\,\mathrm V,3\,\mathrm W$。有电压表 $6\,\mathrm V,20\,\mathrm{k}\Omega$，电流表 $3\,\mathrm A,0.2\,\Omega$ 或 $0.6\,\mathrm A,1\,\Omega$；滑动变阻器 $1000\,\Omega,0.5\,\mathrm A$ 或 $20\,\Omega,2\,\mathrm A$；理想电源 $8\,\mathrm V$。需从零电压起采集12组数据，选择仪器并画连接。
\end{example}
\typeout{ECFIG:fig:c6:example:10:s}\fig{fig:c6:example:10:s}{待测灯泡及给定电源；测量仪器与调压接线由题目要求选择}{\begin{tikzpicture}

\draw(0,0)to[lamp,l={$L:\ 6\,\mathrm V,3\,\mathrm W$}](4,0);
\draw[flow](.3,1.1)--(1.5,1.1)node[midway,above]{$I_L$};
\draw[<->](0,-1)--(4,-1)node[midway,below]{$U_L$};
\draw(6,0)to[battery1,l={$8\,\mathrm V$}](9,0);\node at(7.5,-1){理想电源};

\end{tikzpicture}}
\end{question}
\begin{solution}\textbf{由额定电流和电压表漏流选择仪表}\par\nobreak

图\ref{fig:c6:example:10:s}灯泡额定电流为 $I_N=P_N/U_N=0.50\,\mathrm A$，选 $0.6\,\mathrm A$ 电流表可获得较大偏转。电压表选给定 $6\,\mathrm V$ 表，在额定点漏流仅为 $6/20000=0.00030\,\mathrm A$。测到总电流 $0.50030\,\mathrm A$，仍在 $0.6\,\mathrm A$ 范围内。


\par\medskip\Needspace{4\baselineskip}\noindent\textbf{用分压接法实现从零连续调压}\par\nobreak

\typeout{ECFIG:fig:c6:example:10:p}\fig{fig:c6:example:10:p}{固定端 $A,B$ 跨八伏电源；滑片 $P$ 经电流表向灯泡供电，电压表只跨灯泡}{\begin{tikzpicture}

\draw(0,0)to[battery1,invert,l={$E$}](0,5)--(2,5)--(2,4.5);
\draw(1.8,.5)rectangle(2.2,4.5);\draw(2,.5)--(2,0)--(0,0);
\draw[<-](2.2,2.4)--(3.1,2.4);\draw(3.1,2.4)to[ammeter,l={$R_A$}](5,2.4)--(6.5,2.4)to[lamp,l={$L$}](6.5,0)--(2,0);
\node[left]at(1.7,3.7){$R_P$};\node[above]at(2,5){$A$};\node[below]at(2,0){$B$};\node[above]at(3.1,2.4){$P$};
\draw[->,gray](3.2,1.9)--(3.2,.4)node[midway,right]{起始移到 $B$};
\fill(6.5,2.4)circle(2pt);\fill(6.5,0)circle(2pt);
\draw(6.5,2.4)--(9,2.4)to[voltmeter,l={$R_V$}](9,0)--(6.5,0);
\end{tikzpicture}}

图\ref{fig:c6:example:10:p}选 $20\,\Omega,2\,\mathrm A$ 滑动变阻器，固定端跨电源，滑片作输出。该较小阻值的分压器较易维持负载电压；其允许电流也更充足。接通前 $P$ 移至低电势端 $B$，此时输出对 $B$ 为零。


\par\medskip\Needspace{4\baselineskip}\noindent\textbf{按灯泡实际端电压逐点采集}\par\nobreak

逐渐把 $P$ 移向 $A$，同时读取 $U_V,I_A$，在低压弯曲较明显处适当加密，共取12组。若需扣除有限电压表的分流，每点有
\[
 U_L=U_V,\qquad I_L=I_A-U_V/R_V.
\]
达到额定 $6\,\mathrm V$ 后停止升压。电流表在额定点压降约 $0.5003\,\mathrm V$，电源有足够电压余量。若电压表改跨电流表与灯泡的总支路，就会把这一压降混入灯压。
\end{solution}
\begin{question}\begin{example}[灯泡曲线与串联电阻功率]\label{c6:example:11}
某灯泡伏安曲线在 $U=1,2,4,6,8\,\mathrm V$ 时电流分别约为 $0.20,0.29,0.40,0.46,0.50\,\mathrm A$。将其与定值电阻 $R$ 串联接电源 $E=6.0\,\mathrm V,r=1.0\,\Omega$，实测电流 $0.40\,\mathrm A$。求 $R$ 的功率；另以灯泡各工作点的功率 $P_L$ 作 $P_L-U^2$、$P_L-I^2$ 图，判断哪一关系的相邻采样点割线斜率随工作点增大。
\end{example}
\typeout{ECFIG:fig:c6:example:11:s}\fig{fig:c6:example:11:s}{内阻 $r$、外接电阻 $R$ 与灯泡串联，通过同一电流}{\begin{tikzpicture}

\draw(0,0)to[battery1,invert,l={$E$}](0,3)to[R,l={$r$}](3,3)to[R,l={$R$}](6,3)to[lamp,l={$L$}](6,0)--(0,0);
\draw[flow](.6,4.1)--(1.8,4.1)node[midway,above]{$I$};

\end{tikzpicture}}
\end{question}
\begin{solution}\textbf{先按实测电流读取灯压并列回路方程}\par\nobreak

图\ref{fig:c6:example:11:s}中各串联元件通过同一电流。给定曲线上 $I=0.40\,\mathrm A$ 对应 $U_L=4.0\,\mathrm V$。回路方程 $E=Ir+IR+U_L$ 给出
\[
U_R=6.0-0.40\times1.0-4.0=1.6\,\mathrm V,
\qquad R=4.0\,\Omega,\quad P_R=0.64\,\mathrm W.
\]


\par\medskip\Needspace{4\baselineskip}\noindent\textbf{把灯泡的各采样点转换成功率坐标}\par\nobreak

\typeout{ECFIG:fig:c6:example:11:p}\fig{fig:c6:example:11:p}{同一组灯泡数据的两种变换；连接相邻点的线段用于比较割线斜率}{\begin{tikzpicture}

\begin{scope}
\draw[axis](0,0)--(5.8,0)node[right]{$I^2/\mathrm A^2$};\draw[axis](0,0)--(0,4.7)node[above]{$P_L/\mathrm W$};
\draw[thick]plot coordinates{(.8,.2)(1.682,.58)(3.2,1.6)(4.232,2.76)(5,4)};
\foreach \x/\y in {.8/.2,1.682/.58,3.2/1.6,4.232/2.76,5/4}{\fill(\x,\y)circle(2pt);}
\node[below]at(5,0){$0.25$};\node[left]at(0,4){$4$};
\end{scope}
\begin{scope}[xshift=8cm]
\draw[axis](0,0)--(5.8,0)node[right]{$U^2/\mathrm V^2$};\draw[axis](0,0)--(0,4.7)node[above]{$P_L/\mathrm W$};
\draw[thick]plot coordinates{(.078125,.2)(.3125,.58)(1.25,1.6)(2.8125,2.76)(5,4)};
\foreach \x/\y in {.078125/.2,.3125/.58,1.25/1.6,2.8125/2.76,5/4}{\fill(\x,\y)circle(2pt);}
\node[below]at(5,0){$64$};\node[left]at(0,4){$4$};
\end{scope}

\end{tikzpicture}}

各点 $P_L=UI$，依次为 $0.20,0.58,1.60,2.76,4.00\,\mathrm W$。图\ref{fig:c6:example:11:p}左图横坐标换成 $I^2$，右图换成 $U^2$；没有把恒定的外接 $R$ 的功率与灯泡功率混在一起。


\par\medskip\Needspace{4\baselineskip}\noindent\textbf{直接计算相邻点斜率并区分割线与切线}\par\nobreak

按横坐标为 $I^2$ 的四段分别计算，
\[
\begin{aligned}
\frac{0.58-0.20}{0.29^2-0.20^2}&=\frac{0.38}{0.0441}\simeq8.62\,\Omega,\\
\frac{1.60-0.58}{0.40^2-0.29^2}&=\frac{1.02}{0.0759}\simeq13.44\,\Omega,\\
\frac{2.76-1.60}{0.46^2-0.40^2}&=\frac{1.16}{0.0516}\simeq22.48\,\Omega,\\
\frac{4.00-2.76}{0.50^2-0.46^2}&=\frac{1.24}{0.0384}\simeq32.29\,\Omega.
\end{aligned}
\]
换成横坐标 $U^2$，功率差不变，分母改为电压平方差：
\[
 \begin{aligned}
 \frac{0.38}{2^2-1^2}&=\frac{0.38}{3}\simeq0.1267\,\Omega^{-1},\\
 \frac{1.02}{4^2-2^2}&=\frac{1.02}{12}=0.0850\,\Omega^{-1},\\
 \frac{1.16}{6^2-4^2}&=\frac{1.16}{20}=0.0580\,\Omega^{-1},\\
 \frac{1.24}{8^2-6^2}&=\frac{1.24}{28}\simeq0.0443\,\Omega^{-1}.
 \end{aligned}
\]
因此前者逐渐增大，后者逐渐减小。$P_L/I^2=U/I$ 和 $P_L/U^2=I/U$ 只是过原点割线的比值；仅凭它们增减，不能直接断言任意位置切线斜率的变化。所给离散点支持的是上述采样范围内的趋势。
\end{solution}
\begin{question}\begin{example}[水果电池的仪表改装与拟合]\label{c6:example:12}
表头 $R_g=15\,\Omega,I_g=2\,\mathrm{mA}$，改装为 $4\,\mathrm V$ 表，另有内阻 $2\,\Omega$ 的电流表。水果电池约 $E=4\,\mathrm V,r=200\,\Omega$。数据 $I/\mathrm{mA}=4,5,8,10,12,14$，$U/\mathrm V=3.04,2.85,2.30,1.90,1.50,1.14$。求改装电阻、拟合参数，并选择较小系统误差接法。
\end{example}
\typeout{ECFIG:fig:c6:example:12:s}\fig{fig:c6:example:12:s}{两种候选接法均把电池内阻画出；电流表示数所在分支不同}{\begin{tikzpicture}
\begin{scope}[scale=.85]\draw(0,0)to[battery1,invert,l={$E$}](0,1.7)to[R,l={$r$}](0,4);
\draw[flow](.5,4.8)--(1.6,4.8)node[midway,above]{$I_{\text{源}}$};
\draw(0,4)--(3,4)to[ammeter,l={$R_A$}](3,2)to[R,l={$R_x$}](3,0)--(0,0);
\draw(3,4)--(6,4)to[voltmeter,l={$R_V$}](6,0)--(3,0);
\fill(3,4)circle(2pt);\fill(3,0)circle(2pt);
\draw[flow](2.1,3.5)--(2.1,2.6)node[midway,left]{$I$};
\node at(3,-1){甲：$I_{\text{源}}=I+U/R_V$};
\end{scope}\begin{scope}[xshift=7.5cm,scale=.85]\draw(0,0)to[battery1,invert,l={$E$}](0,1.7)to[R,l={$r$}](0,4);
\draw[flow](.5,4.8)--(1.6,4.8)node[midway,above]{$I_{\text{源}}$};
\draw(0,4)to[ammeter,l={$R_A$}](3,4)--(4.5,4)to[R,l={$R_x$}](4.5,0)--(0,0);
\draw(4.5,4)--(7,4)to[voltmeter,l={$R_V$}](7,0)--(4.5,0);
\fill(4.5,4)circle(2pt)node[above]{$P$};\fill(4.5,0)circle(2pt)node[below]{$Q$};
\draw[flow](.5,3.2)--(1.7,3.2)node[midway,below]{$I$};
\node at(3.5,-1){乙：$I_{\text{源}}=I$，$U=U_{PQ}$};
\end{scope}
\end{tikzpicture}}
\end{question}
\begin{solution}\textbf{先求改装电压表的总内阻}\par\nobreak

满偏电流 $2\,\mathrm{mA}$ 不变，要使整个电压表满量程为 $4\,\mathrm V$，须有
\[
R_V=\frac4{0.002}=2000\,\Omega,\qquad R_s=2000-15=1985\,\Omega.
\]


\par\medskip\Needspace{4\baselineskip}\noindent\textbf{按电池总电流比较两种加载模型}\par\nobreak

图\ref{fig:c6:example:12:s}甲中 $I_{\text{源}}=I+U/R_V$，所以
\[
 U=\frac{ER_V}{R_V+r}-\frac{rR_V}{R_V+r}I.
\]
截距、斜率电阻均比真实值低约 $1-2000/2200=9.1\%$。乙中电流表读总电流，而电压不含电流表压降，故 $U=E-(r+R_A)I$，斜率仅多 $2\,\Omega$，相对约 $1\%$，选乙更合适。

\typeout{ECFIG:fig:c6:example:12:p}\fig{fig:c6:example:12:p}{六组数据与最小二乘直线；横轴用毫安，拟合斜率数值必须相应换算}{\begin{tikzpicture}

\draw[axis](0,0)--(7.6,0)node[right]{$I/\mathrm{mA}$};\draw[axis](0,0)--(0,4.4)node[above]{$U/\mathrm V$};
\draw[thick,domain=0:15]plot({.45*\x},{3.80920-.191041*\x});
\foreach \x/\y in {4/3.04,5/2.85,8/2.30,10/1.90,12/1.50,14/1.14}{\fill(.45*\x,\y)circle(2.5pt);}
\node[left]at(0,3.8092){$3.809$};\node[below]at(6.3,0){$14$};
\node at(4.8,3.7){$U=3.80920-0.191041\,I_{\mathrm{mA}}$};

\end{tikzpicture}}


\par\medskip\Needspace{4\baselineskip}\noindent\textbf{用全部数据求直线参数并恢复真实内阻}\par\nobreak

为避免单位混淆，先记$x=I/(1\,\mathrm{mA})$、$y=U/(1\,\mathrm V)$，拟合$y=A-Bx$。目标函数为$S=\sum_j(y_j-A+Bx_j)^2$。求偏导并令其为零：
\begin{align*}
\frac{\partial S}{\partial A}&=-2\sum_j(y_j-A+Bx_j)=0,\\
\frac{\partial S}{\partial B}&=2\sum_jx_j(y_j-A+Bx_j)=0.
\end{align*}
第一式给$A=\bar y+B\bar x$；代入第二式并整理，得
\[
B=-\frac{\sum x_jy_j-(\sum x_j)(\sum y_j)/6}
{\sum x_j^2-(\sum x_j)^2/6}.
\]
逐项求和：
\[
\sum x_j=53,\quad \sum y_j=12.73,\quad
\sum x_j^2=545,\quad \sum x_jy_j=97.77.
\]
所以$\bar x=53/6$、$\bar y=12.73/6$，并有
\[
B=-\frac{97.77-53\times12.73/6}{545-53^2/6}
=0.191041,\qquad
A=\frac{12.73}{6}+B\frac{53}{6}=3.80920.
\]
图\ref{fig:c6:example:12:p}的带符号斜率为 $-0.191041\,\mathrm{V/mA}$。因为 $x=1000I/(1\,\mathrm A)$，恢复电压、电流单位后，
\[
 U\simeq3.80920\,\mathrm V-(1000B\,\Omega)I
 =3.80920\,\mathrm V-(191.041\,\Omega)I.
\]
与乙接法的 $U=E-(r+R_A)I$ 比较，
\[
 E=A\,\mathrm V\simeq3.81\,\mathrm V,\qquad
 r=1000B\,\Omega-R_A=(191.041-2)\,\Omega\simeq189\,\Omega.
\]
横轴若用毫安，斜率绝对值须乘1000才得到欧姆。计算中间位保留较多，最终精度仍受给定读数限制。
\end{solution}
\begin{question}\begin{example}[不拆表头的测阻与量程标定]\label{c6:example:13}
表头满偏电流 $I_g$ 已知、内阻 $R_g$ 待测，拟加串联电阻箱 $R_s$ 改电压表。标准电压表跨分压器输出与回路公共端，双掷开关挡1使输出直接接表头，挡2使输出经 $R_s$ 后接表头。说明怎样在不拆表头的条件下测 $R_g$ 并标定新量程。
\end{example}
\typeout{ECFIG:fig:c6:example:13:s}\fig{fig:c6:example:13:s}{双掷开关公共端接滑片输出；标准表跨整个输出端口，接点用实心点表示}{\begin{tikzpicture}

\draw(0,0)to[battery1,invert,l={$E$}](0,5.8)--(2,5.8)--(2,4.8);
\draw(1.8,.6)rectangle(2.2,4.8);\draw(2,.6)--(2,0)--(0,0);\draw(2,0)--(10,0);
\node[left]at(1.7,4.1){$R_P$};\draw[<-](2.2,3)--(3.6,3);\draw(3.6,3)--(4.5,3);
\fill(3.6,3)circle(2pt)node[below]{$P$};\fill(4.5,3)circle(2pt);
\fill(5.5,3.8)circle(2pt)node[above]{1};\fill(5.5,2.2)circle(2pt)node[below]{2};
\draw[thick](4.5,3)--(5.3,3.65);
\draw(5.5,3.8)--(8,3.8)--(8,1.8)to[ammeter,l={$G$}](8,0);
\draw(5.5,2.2)to[R,l={$R_s$}](8,2.2);\fill(8,2.2)circle(2pt);
\draw(3.6,3)--(3.6,5.2)--(10,5.2)to[voltmeter,l={$V_{\text{标}}$}](10,0);
\fill(8,0)circle(2pt)node[below]{$Q$};\fill(10,0)circle(2pt);

\end{tikzpicture}}
\end{question}
\begin{solution}\textbf{在挡一测裸表头满偏电压}\par\nobreak

图\ref{fig:c6:example:13:s}标准表始终跨输出 $P,Q$。输出先调至最小，开关接1，再逐渐升高到表头满偏。此时 $R_s$ 支路断开，标准表直接读取 $G$ 的实际电压 $U_g$，所以 $R_g=U_g/I_g$。

\typeout{ECFIG:fig:c6:example:13:p}\fig{fig:c6:example:13:p}{两种开关状态的实际测量端口；挡二标准表跨串联电阻与表头的整体}{\begin{tikzpicture}

\begin{scope}
\draw(0,3)--(3,3)to[ammeter,l={$R_g$}](3,0)--(0,0);
\draw(0,3)--(-2,3)to[voltmeter,l={$U_g$}](-2,0)--(0,0);
\fill(0,3)circle(2pt)node[above]{$P$};\fill(0,0)circle(2pt)node[below]{$Q$};
\draw[flow](3.9,2.7)--(3.9,1.8)node[midway,right]{$I_g$};\node at(1,-1){挡1：直接接表头};
\end{scope}
\begin{scope}[xshift=8.8cm]
\draw(0,3)--(3,3)to[R,l={$R_s$}](3,1.6)to[ammeter,l={$R_g$}](3,0)--(0,0);
\draw(0,3)--(-2,3)to[voltmeter,l={$U_m$}](-2,0)--(0,0);
\fill(0,3)circle(2pt)node[above]{$P$};\fill(0,0)circle(2pt)node[below]{$Q$};
\draw[flow](4.2,2.7)--(4.2,1.8)node[midway,right]{$I_g$};\node at(1,-1){挡2：串入电阻箱};
\end{scope}

\end{tikzpicture}}


\par\medskip\Needspace{4\baselineskip}\noindent\textbf{在挡二计算串联阻值并核对目标量程}\par\nobreak

先降低输出，再换到2。图\ref{fig:c6:example:13:p}右图中表头与 $R_s$ 串联，若目标量程为 $U_m$，设
\[
R_s=\frac{U_m}{I_g}-R_g.
\]
由低电压逐渐升高，标准表到 $U_m$ 时检查表头恰满偏；再在半量程等若干点对照。不能在裸表头挡直接施加整个新量程。


\par\medskip\Needspace{4\baselineskip}\noindent\textbf{说明标准表加载为何不改变测阻关系}\par\nobreak

标准表有限内阻会增加分压器的负载，但它测到的正是同一对 $P,Q$ 端点的实际电压。因此 $U_g/I_g$ 不依赖分压器的空载输出值。挡二右支路的 $I_g$ 仍是表头电流，不能误用包含标准表分流的分压器总输出电流。全过程不需要拆开表头接线。
\end{solution}
